\chapter{Deployment, DevOps, and Operations}

\section{Environment Topology}
A practical environment model includes:
\begin{itemize}[leftmargin=1.2cm]
  \item Development: local monorepo with backend and frontend concurrently.
  \item Staging: cloud VM or container stack with test data.
  \item Production: hardened setup with managed database and observability stack.
\end{itemize}

\section{Configuration Management}
Environment variables control ports, database URI, JWT secret, OAuth credentials, frontend origin, and access codes. Secret values must never be committed to source control.

\section{Build and Run Pipeline}
\subsection{Root Orchestration}
The monorepo root starts backend and frontend together for integrated development testing.

\subsection{Backend Runtime}
Backend service starts HTTP server, websocket endpoint, MQTT broker, and DB connection in one process.

\subsection{Frontend Runtime}
Vite serves SPA with hot-module reload in development and optimized static build in production.

\section{Containerization Strategy (Recommended)}
A production-ready containerization path:
\begin{enumerate}[leftmargin=1.2cm]
  \item Backend image with Node runtime, non-root user, healthcheck endpoint.
  \item Frontend static build served by Nginx or edge hosting.
  \item Externalized MongoDB as managed database service.
\end{enumerate}

\section{CI/CD Blueprint}
\begin{itemize}[leftmargin=1.2cm]
  \item Trigger on pull request and main branch merge.
  \item Install dependencies and run lint checks.
  \item Run automated test suites.
  \item Build artifacts and publish deploy package.
  \item Deploy to staging, then production with approval gate.
\end{itemize}

\section{Observability Architecture}
\subsection{Logs}
Adopt structured logs with request IDs and event correlation IDs.

\subsection{Metrics}
Capture HTTP latency, broker throughput, socket connection counts, and emergency incidents.

\subsection{Alerts}
Configure alerting for service downtime, DB connectivity issues, and abnormal emergency burst patterns.

\section{Backup and Recovery}
\begin{itemize}[leftmargin=1.2cm]
  \item Daily MongoDB backups with retention policy.
  \item Configuration backup for environment templates.
  \item Recovery drill procedure tested periodically.
\end{itemize}

\section{Operational Runbook Essentials}
Operational teams should maintain runbooks for:
\begin{itemize}[leftmargin=1.2cm]
  \item Incident triage and escalation.
  \item Broker connection troubleshooting.
  \item Emergency false alert investigation.
  \item User account recovery and access revocation.
\end{itemize}

\section{Scalability Roadmap}
\begin{itemize}[leftmargin=1.2cm]
  \item Separate broker and API services.
  \item Introduce Redis adapter for horizontal socket scaling.
  \item Add telemetry queueing for burst management.
  \item Partition data and house-level workloads.
\end{itemize}

\section{Operational KPIs}
\begin{table}[H]
\centering
\caption{Sample operational service-level indicators}
\begin{tabular}{p{5cm}p{3cm}p{6cm}}
\toprule
Indicator & Target & Notes \\
\midrule
API Availability & > 99.5\% & Monthly rolling target \\
p95 API Latency & < 350 ms & Authenticated business routes \\
Socket Delivery Success & > 99\% & Event emission to intended rooms \\
Emergency Processing Delay & < 2 s & Sensor receipt to user alert \\
\bottomrule
\end{tabular}
\end{table}

\section{Operations Summary}
The project is deployment-ready for educational and pilot usage. With incremental hardening in CI/CD, observability, and scaling controls, it can evolve into production-grade service operations.
